\section{The \$200 billion machine that forgot to cure anything}\label{19_200_billion-the-200-billion-machine-that-forgot-to-cure-anything}

Biomedical research dysfunction is not a methamphetamine-specific anomaly --- it is a system-wide structural failure spanning every major disease category, driven by misaligned incentives, and insulated from reform by one of the most effective institutional capture arrangements in modern governance. \textbf{Only 5\% of animal-tested therapies ever achieve regulatory approval}, according to the most comprehensive meta-analysis to date. The pattern your document identifies --- mouse studies that don't translate, self-dealing peer review, career incentives rewarding publications over clinical outcomes --- is the norm across Alzheimer's, cancer, depression, addiction, and virtually every other field. Despite documented fraud consuming billions in public funds and a reproducibility crisis acknowledged by \textbf{72\% of researchers themselves}, \href{https://www.barnesandnoble.com/w/rigor-mortis-richard-f-harris/1124804938}{Barnes \& Noble +2} there has been no prosecution, no Congressional investigation, no structural reform. The reason is not a mystery: the current system generates \$200+ billion annually for universities, publishers, equipment manufacturers, and career scientists --- a concentrated beneficiary class that vastly outweighs the diffuse public interest in cures.

\subsection{The translation failure is universal and quantified}\label{19_200_billion-the-translation-failure-is-universal-and-quantified}

The foundational claim of mouse-model biomedical research --- that animal studies reliably predict human outcomes --- has been systematically disproven across decades of evidence. A landmark 2024 umbrella review in \emph{PLOS Biology} by Ineichen et al.~analyzed 122 systematic reviews covering 367 therapeutic interventions across 54 diseases \href{https://journals.plos.org/plosbiology/article?id=10.1371/journal.pbio.3002667}{PLOS} \href{https://pubmed.ncbi.nlm.nih.gov/38870090/}{PubMed} and found that while 50\% of animal-tested therapies progressed to some human study, \textbf{only 5\% achieved regulatory approval}. \href{https://www.eurekalert.org/news-releases/1046892}{EurekAlert!} The median timeline from animal result to approval stretches to 10 years, with neurological diseases accounting for a disproportionate 32\% of the failed attempts. \href{https://journals.plos.org/plosbiology/article?id=10.1371/journal.pbio.3002667}{PLOS} \href{https://www.eurekalert.org/news-releases/1046892}{EurekAlert!}

This is not new knowledge. Hackam and Redelmeier reported in \emph{JAMA} in 2006 that only about one-third of highly cited animal research translates even to the level of human randomized trials. \href{https://journals.plos.org/plosmedicine/article?id=10.1371/journal.pmed.1000245}{PLOS} \href{https://link.springer.com/article/10.1186/s12967-019-1976-2}{Springer} Contopoulos-Ioannidis found in 2003 that only 5\% of ``high-impact'' basic science discoveries claiming clinical relevance produced approved agents within a decade. \href{https://www.droracle.ai/articles/236934/what-percentage-of-promising-animal-studies-fail-to-translate}{Drooracle} Pound et al.~asked in the \emph{BMJ} in 2004, ``Where is the evidence that animal research benefits humans?'' --- and found remarkably little.

The failure is not abstract. In methamphetamine addiction, \textbf{zero FDA-approved medications exist} \href{https://oasas.ny.gov/methamphetamine}{Office of Addiction Services and Supports} despite decades of mouse research on dopaminergic pathways. The FDA itself acknowledged in October 2023 that ``currently there is no FDA-approved medication for stimulant use disorder.'' \href{https://www.statnews.com/2023/10/04/methamphetamine-cocaine-addiction-treatment-fda/}{Statnews} In Alzheimer's, the drug failure rate reached \textbf{99.6\%} between 2002 and 2012 --- 244 compounds tested, one approved. \href{https://www.scientificamerican.com/article/why-alzheimer-s-drugs-keep-failing/}{Scientific American} In cancer, translation rates from animal models run below 8\%, \href{https://pmc.ncbi.nlm.nih.gov/articles/PMC3902221/}{PubMed Central} with oncology drugs showing the lowest overall FDA approval rate at 3.4\%. The serotonin hypothesis of depression, which justified decades of SSRI prescriptions, was demolished by Moncrieff et al.'s 2022 umbrella review in \emph{Molecular Psychiatry}, \href{https://www.ucl.ac.uk/news/2022/jul/no-evidence-depression-caused-low-serotonin-levels-finds-comprehensive-review}{UCL} which found ``no consistent evidence'' linking serotonin levels to depression across every major research methodology. \href{https://psychopharmacologyinstitute.com/section/the-serotonin-theory-of-depression-a-systematic-umbrella-review-of-the-evidence-2683-5270/}{Psychopharmacology Institute}

The Amgen replication study crystallized the crisis: scientists attempted to replicate 53 ``landmark'' preclinical cancer studies and \textbf{succeeded with only 6 --- an 11\% replication rate}. \href{https://www.kirkusreviews.com/book-reviews/richard-harris/rigor-mortis-sloppy-science/}{Kirkus Reviews} \href{https://issues.org/ending-reproducibility-crisis-medical-research-brownlee-bielekova/}{Issues} Bayer Healthcare replicated only 36\% of studies it tested. \href{https://www.castoredc.com/blog/replication-crisis-in-clinical-research/}{Castor} Glasziou and Chalmers estimated that \textbf{85\% of biomedical research funding is avoidably wasted} on studies with poor methodology and reporting. \href{https://www.aje.com/arc/pre-registration-vs-registered-reports}{AJE} \href{https://issues.org/ending-reproducibility-crisis-medical-research-brownlee-bielekova/}{Issues} This is not a fringe claim --- it was published in \emph{The Lancet}.

\subsection{How the incentive machine actually works}\label{19_200_billion-how-the-incentive-machine-actually-works}

The system producing these outcomes is not broken in the conventional sense. It functions precisely as designed --- for the people it was designed to serve. Understanding this requires examining the concrete mechanics of how money, careers, and prestige flow through the biomedical research enterprise.

NIH distributes approximately \textbf{\$47--48 billion annually}, \href{https://www.endocrine.org/advocacy/position-statements/biomedical-research-funding}{Endocrine Society} with 82\% going to extramural grants \href{https://www.gao.gov/assets/gao-23-105656.pdf}{U.S. GAO} supporting 300,000+ researchers at 2,500+ institutions. \href{https://www.nih.gov/about-nih/organization/budget}{National Institutes of Health} The core funding mechanism is the R01 grant: up to \$250,000 per year in direct costs for five years. \href{https://www.niaid.nih.gov/grants-contracts/research-project-grants}{NIAID} But universities add \textbf{indirect cost rates of 50--60\%} on top of direct costs, \href{https://www.bdo.com/insights/industries/nonprofit-education/understanding-the-nih-s-new-indirect-cost-rate-policy-what-nonprofit-and-higher-education-cfos-need}{BDO} meaning a \$250,000 grant actually costs NIH \$375,000--400,000, with the difference going to institutional overhead. Total overhead payments to universities exceed \textbf{\$9 billion annually}. \href{https://www.socialsciencespace.com/2025/02/those-indirect-costs-target-by-doge-directly-support-americas-research-excellence/}{Social Science Space} \href{https://www.science.org/content/article/nih-slashes-overhead-payments-research-sparking-outrage}{Science} This overhead funds building debt service, administrative salaries, and institutional infrastructure --- creating a powerful financial incentive for universities to maximize grant volume regardless of research quality. \href{https://pmc.ncbi.nlm.nih.gov/articles/PMC4000813/}{PubMed Central}

The R01 structure inherently favors mouse studies because a typical mouse experiment costs \textbf{\$100,000--200,000} and fits neatly within a five-year grant timeline that matches mouse lifespan. A Phase II human clinical trial costs \textbf{\$7--20 million} \href{https://www.abacum.ai/blog/clinical-trial-costing}{Abacum} --- 50 to 100 times more --- and requires infrastructure, regulatory expertise, and timelines that no single R01 can support. The system doesn't merely prefer mouse studies; it is architecturally designed to produce them.

Who decides which grants get funded? ``Study sections'' composed of the same academic researchers who apply for grants. \href{https://grants.nih.gov/grants-process/review/first-level}{NIH Grants} NIH prohibits direct financial conflicts, \href{https://pmc.ncbi.nlm.nih.gov/articles/PMC5825276/}{PubMed Central} \href{https://grants.nih.gov/grants-process/review/first-level}{NIH Grants} but the structural bias is inescapable: mouse researchers evaluate mouse proposals. \textbf{95\% of institutions worldwide} require peer-reviewed publications for tenure, and \textbf{67\% require grant funding}. \href{https://pmc.ncbi.nlm.nih.gov/articles/PMC7315647/}{PubMed Central} \href{https://academyhealth.org/sites/default/files/publication/\%5Bfield_date:custom:Y\%5D-\%5Bfield_date:custom:m\%5D/academicincentivesresearchimpact_feb2021.pdf}{AcademyHealth} A researcher's career survival depends on publishing prolifically in the paradigm that peer reviewers --- working in that same paradigm --- will fund. The average age of a first independent NIH grant is now \textbf{42 for PhD holders}, \href{https://www.ncbi.nlm.nih.gov/books/NBK22694/}{NCBI} meaning scientists spend their entire formative careers learning to produce what the system rewards rather than what patients need.

The result is a hypercompetitive environment where, \href{https://www.thetransmitter.org/spectrum/book-review-rigor-mortis-reveals-rampant-sloppiness-science/}{The Transmitter} as a 2024 survey of 1,630 biomedical researchers found, \textbf{62\% said ``pressure to publish'' always or very often contributes to irreproducible research}. \href{https://www.nature.com/articles/d41586-024-04253-w}{Nature +2} Grant success rates have plummeted --- falling to \textbf{18.5--19\% for early-stage investigators in FY2025}, down from nearly 30\% just two years earlier. \href{https://www.statnews.com/2025/12/18/nih-early-career-researchers-grant-success-rate-falls/}{Statnews} In this environment, the rational career strategy is to generate publishable results fast and cheap (mice), not to conduct expensive, slow, uncertain human studies that might produce null results fatal to a tenure case. \href{https://www.thetransmitter.org/spectrum/book-review-rigor-mortis-reveals-rampant-sloppiness-science/}{The Transmitter}

\subsection{The Alzheimer's scandal proves the system cannot self-correct}\label{19_200_billion-the-alzheimers-scandal-proves-the-system-cannot-self-correct}

The 2022 Alzheimer's amyloid fraud case is the single most important test of whether the biomedical research system has functioning accountability mechanisms. It failed every test.

In July 2022, \emph{Science} journalist Charles Piller published an investigation revealing that Sylvain Lesné, a neuroscientist at the University of Minnesota, had apparently manipulated Western blot images \href{https://www.thetransmitter.org/science-and-society/doctored-fraud-arrogance-and-tragedy-in-the-quest-to-cure-alzheimers-an-excerpt/}{The Transmitter} in his landmark 2006 \emph{Nature} paper \href{https://healthjournalism.org/blog/2022/09/how-i-did-it-reporting-on-apparent-misconduct-in-influential-alzheimers-research/}{Association of Health Care Journalists} identifying Aβ*56 as a causal agent in Alzheimer's memory loss. \href{https://www.science.org/content/article/potential-fabrication-research-images-threatens-key-theory-alzheimers-disease}{Science} The whistleblower, Vanderbilt neuroscientist Matthew Schrag, \href{https://www.barnesandnoble.com/w/doctored-charles-piller/1145681914}{Barnes \& Noble} identified at least \textbf{70 instances of possible image tampering across 20+ papers}. The 2006 paper had been cited approximately \textbf{2,300 times} and became the fourth most-cited basic-research Alzheimer's paper ever published. \href{https://healthjournalism.org/blog/2022/09/how-i-did-it-reporting-on-apparent-misconduct-in-influential-alzheimers-research/}{Association of Health Care Journalists}

What followed was institutional paralysis. \emph{Nature} issued a cautious ``publisher's note'' and \textbf{did not retract the paper until June 2024 --- two full years later}. \href{https://www.clubvita.net/us/news-and-insights/alzheimers-research-scandal}{Club Vita} The University of Minnesota conducted a 2.5-year investigation \href{https://minnesotareformer.com/briefs/alzheimers-researcher-alleged-to-have-doctored-images-is-leaving-umn/}{Minnesota Reformer} \href{https://www.science.org/content/article/alzheimer-s-scientist-resigns-after-university-finds-data-integrity-concerns-papers}{Science} that initially ``cleared'' Lesné of misconduct \href{https://minnesotareformer.com/briefs/alzheimers-researcher-alleged-to-have-doctored-images-is-leaving-umn/}{Minnesota Reformer} on the \emph{Nature} paper before eventually recommending retraction of four additional papers. Lesné quietly resigned effective March 2025 \href{https://www.science.org/content/article/alzheimer-s-scientist-resigns-after-university-finds-data-integrity-concerns-papers}{Science} --- never fired, never prosecuted, never publicly sanctioned. \href{https://www.theburningplatform.com/2026/03/26/the-great-alzheimers-scam-and-the-proven-cures-theyve-buried-for-billions-2/}{The Burning Platform} Four months after Schrag reported the fraud to NIH, \textbf{Lesné was awarded a new \$765,000 NIH grant}, with the program officer being a co-author on the very paper under investigation. \href{https://www.clubvita.net/us/news-and-insights/alzheimers-research-scandal}{Club Vita}

The downstream consequences were nil. \textbf{No Congressional hearing} was held on the fraud. \textbf{No criminal prosecution} was initiated despite potential False Claims Act violations on approximately \$8 million in grants. NIH did not change its Alzheimer's funding priorities --- it continued spending approximately \textbf{\$1.6 billion annually on amyloid-related projects}, roughly half of all federal Alzheimer's funding. \href{https://articles.mercola.com/sites/articles/archive/2026/03/26/alzheimers-amyloid-hypothesis-fraud.aspx}{Mercola} \href{https://katiecouric.com/health/alzheimers-research-fraud-matthew-schrag-sylvain-lesne/}{Katie Couric} No GAO investigation examined whether decades of amyloid-focused funding was effectively allocated. The university \textbf{promoted the administrator in charge of research integrity} during the scandal period. \href{https://www.startribune.com/umn-professors-fabrication-plagiarism-scandals/601554124}{Star Tribune}

The fraud's downstream impact is staggering. About \textbf{100 of 130 Alzheimer's drugs in clinical trials} at the time targeted amyloids. \href{https://katiecouric.com/health/alzheimers-research-fraud-matthew-schrag-sylvain-lesne/}{Katie Couric} Total NIH Alzheimer's funding grew from \$448 million to approximately \$3.9 billion annually \href{https://www.sgtreport.com/2026/03/the-great-alzheimers-scam-and-the-proven-cures-theyve-buried-for-billions-3/}{SGT Report} \href{https://alzimpact.org/research}{Alzimpact} --- a seven-fold increase \href{https://www.alz.org/get-involved-now/advocate/research-funding}{Alzheimer's Association} --- with cumulative spending on amyloid-related research reaching into the tens of billions. Private industry spent additional tens of billions on amyloid-targeting clinical trials. The two anti-amyloid drugs eventually approved (lecanemab and donanemab) show only \textbf{27--30\% slowing of cognitive decline} \href{https://www.michiganmedicine.org/health-lab/everything-you-know-about-alzheimers-wrong}{Michigan Medicine} with serious side effects including brain bleeding.

Piller's 2025 book \emph{Doctored} documented a pattern extending far beyond Lesné. He identified 132 papers by Eliezer Masliah, former director of the Division of Neuroscience at the National Institute on Aging, that were ``riddled with apparently doctored data.'' \href{https://link.springer.com/article/10.1007/s11357-025-01661-2}{Springer} \href{https://www.clubvita.net/us/news-and-insights/alzheimers-research-scandal}{Club Vita} Former Stanford President Marc Tessier-Lavigne stepped down over related concerns. The pattern is systemic, not individual.

\subsection{Why accountability institutions all fail simultaneously}\label{19_200_billion-why-accountability-institutions-all-fail-simultaneously}

The absence of response to documented fraud is not mysterious when analyzed through institutional economics. Every entity that should provide accountability has concrete financial or political reasons not to.

\textbf{Universities} receive \$9+ billion annually in overhead from NIH grants. \href{https://www.socialsciencespace.com/2025/02/those-indirect-costs-target-by-doge-directly-support-americas-research-excellence/}{Social Science Space} Investigating a researcher means potentially losing their grants and the associated overhead --- typically 50--60 cents for every grant dollar. The direct cost of handling a single research misconduct case averages \textbf{\$525,000}, with over \textbf{\$110 million} incurred annually across all institutions. \href{https://pmc.ncbi.nlm.nih.gov/articles/PMC5206685/}{PubMed Central} The financial incentive is overwhelmingly to look the other way.

\textbf{The Office of Research Integrity}, nominally responsible for policing fraud, operates with a budget of approximately \textbf{\$15 million and 37 staff} \href{https://retractionwatch.com/wp-content/uploads/2024/03/ORI-Budget-Grew-to.pdf}{Retraction Watch} to oversee \$47.5 billion in funding and 300,000+ researchers. This represents 0.03\% of the NIH budget. In FY2022, ORI received 269 allegations, closed 42 cases, and made \textbf{only 9 findings of misconduct}. \href{https://cen.acs.org/research-integrity/US-Office-Research-Integrity-received-269-allegations-research-misconduct-last-fiscal-year/101/web/2023/02}{C\&EN} It lacks subpoena power. It has had four directors in ten years. The Lesné fraud was not detected by ORI --- it was exposed by an independent journalist and a junior whistleblower scientist.

\textbf{Scientific journals} profit from publications, not retractions. Elsevier reported profit margins of \textbf{37.8\%} on revenue of €3.26 billion --- margins rivaling Apple and Google --- by publishing content researchers provide for free and selling it back to the universities that employed them. \href{https://pmc.ncbi.nlm.nih.gov/articles/PMC6740196/}{PubMed Central} The ``Big Five'' publishers control approximately 50\% of all global research output. \href{https://theconversation.com/academic-publishing-is-a-multibillion-dollar-industry-its-not-always-good-for-science-250056}{The Conversation} Retracting papers reduces citation metrics and journal prestige; the incentive structure actively discourages quality control. Over \textbf{10,000 papers were retracted globally in 2023}, a record growing at 23\% annually, suggesting the problem is accelerating faster than corrections. \href{https://minnesotareformer.com/briefs/alzheimers-researcher-alleged-to-have-doctored-images-is-leaving-umn/}{Minnesota Reformer}

\textbf{Congress} faces a classic Olsonian collective action problem. NIH funding supports \textbf{390,000+ jobs} across all 50 states and generates an estimated \$94 billion in economic activity. \href{https://www.citizensforethics.org/reports-investigations/crew-reports/doges-big-illusion-the-heavy-costs-of-the-trump-administrations-so-called-efficiency/}{Citizens for Responsibility and Ethics in Washington} \href{https://www.endocrine.org/advocacy/position-statements/biomedical-research-funding}{Endocrine Society} Every member of Congress has NIH-dependent constituents --- research universities, medical centers, biotech companies. The political incentive is to announce funding increases, not to question whether the money produces results. Cutting or restructuring NIH means cutting jobs in your district.

\textbf{Patient advocacy groups}, which should represent the victims, have been captured by pharmaceutical industry funding. The Alzheimer's Association received \textbf{\$4.9 million from pharmaceutical and biotech companies} in FY2023, \href{https://www.j-alz.com/content/time-look-upstream-and-hold-us-alzheimers-association-accountable-its-actions}{Journal of Alzheimer's Disease} hosts corporate sponsorships of \$100,000--\$250,000 at its annual conference, \href{https://www.healthnewsreview.org/2016/03/investigation-alzheimers-association-fighting-didnt-address-role-corporate-sponsorships/}{Healthnewsreview} and actively supported FDA approval of the controversial anti-amyloid drug Aduhelm. \href{https://www.j-alz.com/content/time-look-upstream-and-hold-us-alzheimers-association-accountable-its-actions}{Journal of Alzheimer's Disease} Rather than demanding accountability for research fraud, the Association published forceful rebuttals to those who questioned the amyloid paradigm, calling criticism ``harmful myths.'' \href{https://www.alz.org/news/2025/alzheimers-research-harmful-myths-must-stop}{Alzheimer's Association}

\textbf{The media} cannot sustain coverage because the technical complexity (Western blot manipulation, amyloid oligomers) is inaccessible to general audiences, the connection between a specific fraud and patient harm is indirect and deniable, and pharmaceutical advertising revenue creates financial dependencies. Coverage of the Lesné scandal was reactive --- initial reports followed by silence.

The political science literature maps this precisely. George Stigler's 1971 theory of regulatory capture --- ``regulation is often acquired by the industry and is designed and operated primarily for its benefit'' --- applies directly. Mancur Olson's logic of collective action explains why well-organized small groups (research universities, professional societies, publishers) with high per-member stakes consistently outlobby a diffuse public with small per-person stakes. Edwards and Roy's 2017 analysis in \emph{Environmental Engineering Science} documented how the incentive system ``selectively weeds out ethical and altruistic actors, while selecting for academics who are more comfortable and responsive to perverse incentives.'' \href{https://pmc.ncbi.nlm.nih.gov/articles/PMC8247552/}{PubMed Central}

\subsection{The problem is global because the U.S. exports the incentive structure}\label{19_200_billion-the-problem-is-global-because-the-u.s.-exports-the-incentive-structure}

This dysfunction is not an American peculiarity. NIH is the world's largest biomedical research funder, \href{https://en.wikipedia.org/wiki/Biomedical_research_in_the_United_States}{Wikipedia} and its priorities shape what the global scientific community studies. Researchers worldwide must publish in American journals --- \emph{Cell}, \emph{Science}, \emph{Nature}, \emph{NEJM} --- for career advancement, creating de facto U.S. control over what research gets recognition. The ``publish or perish'' culture has been documented in Europe, Asia, Africa, and developing nations; \href{https://www.kirkusreviews.com/book-reviews/richard-harris/rigor-mortis-sloppy-science/}{Kirkus Reviews} a 2024 survey found that \textbf{95\% of biomedical faculties worldwide} use traditional publication metrics for evaluation. \href{https://pmc.ncbi.nlm.nih.gov/articles/PMC7315647/}{PubMed Central} \href{https://academyhealth.org/sites/default/files/publication/\%5Bfield_date:custom:Y\%5D-\%5Bfield_date:custom:m\%5D/academicincentivesresearchimpact_feb2021.pdf}{AcademyHealth}

The European Research Council distributes approximately €2.3 billion annually \href{https://erc.europa.eu/about-erc/erc-glance}{European Research Council} --- roughly 5\% of NIH's budget --- with similar peer review structures. Germany's DFG, the UK's MRC, France's INSERM, and Japan's JSPS all operate within the same paradigm. China has become the world's second-largest research producer (23.4\% of global output) but suffers even worse fraud: approximately \textbf{three-quarters of 14,000 global retractions in 2023} involved Chinese authors, \href{https://www.chemistryworld.com/news/china-conducts-nationwide-audit-of-research-misconduct-after-thousands-of-papers-retracted/4018981.article}{Chemistry World} driven by cash bonuses for publications and promotion quotas tied to publication counts. \href{https://www.spectator.co.uk/article/the-paper-mills-helping-china-commit-scientific-fraud/}{The Spectator} China's pharmaceutical innovation conversion rate runs below \textbf{8\% annually} versus 50--70\% in developed countries. \href{https://pmc.ncbi.nlm.nih.gov/articles/PMC11974572/}{PubMed Central}

No country has solved translational failure. Overall drug approval success rates run approximately \textbf{12.8\%} globally for candidates from the U.S., EU, or Japan. \href{https://pmc.ncbi.nlm.nih.gov/articles/PMC6409418/}{PubMed Central} Oncology success rates have been \textbf{trending down since 2015}, \href{https://www.science.org/content/blog-post/latest-drug-failure-and-approval-rates}{Science} falling to 5.2\% in 2021. Twenty-five of 27 EU countries have longer median delays to drug availability than the U.S. Japan faces persistent ``drug lag.'' The structural problems transcend borders because they are rooted in the fundamental incentive architecture of modern science, which the U.S. built and exported.

\subsection{Reforms exist on paper but lack political vehicles}\label{19_200_billion-reforms-exist-on-paper-but-lack-political-vehicles}

A robust ecosystem of reform proposals exists. The HHMI ``fund people not projects'' model, providing approximately \$9--11 million over seven-year terms to 320 investigators, produces \textbf{more high-impact papers} than comparably qualified NIH-funded scientists. \href{https://www.journals.uchicago.edu/doi/full/10.1086/699933}{University of Chicago Press} Patrick Collison's Fast Grants distributed \$50+ million during COVID with 48-hour turnaround, \href{https://www.lewistownsentinel.com/news/business/2023/07/grants-from-tech-billionaire-to-speed-up-science-research/}{The Sentinel} and \textbf{64\% of recipients said their funded work wouldn't have happened otherwise}. \href{https://johnhcochrane.blogspot.com/2021/06/fast-grants-and-economics-of.html}{Blogger} The Arc Institute pledged \$650+ million for 8-year no-strings-attached research funding. \href{https://thehill.com/homenews/ap/ap-business/ap-quick-grants-from-tech-billionaires-aim-to-speed-up-science-research-but-not-all-scientists-approve/}{The Hill} ARPA-H, established in 2022 with a \$1.5 billion budget, \href{https://www.science.org/content/article/final-nih-budget-2024-essentially-flat}{Science} adopted the DARPA model of empowered program managers rather than peer review committees. \href{https://www.researchgate.net/publication/337776171_15_Does_NIH_need_a_DARPA}{ResearchGate} Focused Research Organizations --- nonprofit, startup-like entities tackling specific bottleneck problems --- have launched 10 projects since 2021 through Convergent Research. \href{https://bcs.mit.edu/news/former-mit-researchers-advance-new-model-innovation}{Mit}

The Registered Reports movement, now adopted by \textbf{300+ journals}, pre-reviews study designs before data collection, \href{https://spectator.us/bad-science-coronavirus/}{Spectator USA} and evidence shows it dramatically reduces publication bias: 61\% of hypotheses tested in Registered Reports are not supported, versus the near-universal positive results in traditional publishing. \href{https://www.sciencedirect.com/topics/psychology/preregistered}{ScienceDirect} \href{https://en.wikipedia.org/wiki/Replication_crisis}{Wikipedia} The Declaration on Research Assessment (DORA), signed by thousands of organizations, calls for evaluating research on its merits rather than journal impact factors. \href{https://pmc.ncbi.nlm.nih.gov/articles/PMC5206685/}{PubMed Central} Stuart Buck's Good Science Project published a 35-page NIH reform plan. \href{https://goodscience.substack.com/p/new-in-metascience-and-good-government}{Substack} John Ioannidis has proposed everything from stricter statistical thresholds to a complete overhaul of researcher evaluation systems. \href{http://backreaction.blogspot.com/2018/11/book-review-rigor-mortis-by-richard.html}{Blogger}

Legislatively, the picture is thin. Senator Bill Cassidy released an NIH modernization white paper in 2024. \href{https://goodscienceproject.org/articles/metascience-reforms-at-nih-2/}{Goodscienceproject} The House Energy and Commerce Committee proposed consolidating NIH's 27 institutes into 15, \href{https://www.apaservices.org/advocacy/news/opposing-nih-radical-reorganization}{APA Services} \href{https://ww2.aip.org/fyi/republicans-push-nih-reform}{AIP} mentioning research misconduct as one issue among many. \href{https://www.help.senate.gov/imo/media/doc/nih_modernization_5924pdf.pdf}{U.S. Senate Committee on Health, Education, Labor \& Pensions} Congress included language in FY2025 appropriations directing NIH to spend \textbf{at least \$10 million on replication experiments} \href{https://www.timeshighereducation.com/news/congress-pushes-nih-review-reproducibility-and-near-misses}{Times Higher Education} \href{https://goodscience.substack.com/p/metascience-reforms-at-nih}{Substack} --- a rounding error in a \$48 billion budget. The False Claims Act has been used to recover funds from universities (Duke paid \textbf{\$112.5 million} in a single settlement), \href{https://www.wiley.law/newsletter-Got-Grants-Public-andPrivate-Sector-Federal-Grant-Recipients-Not-Immune-to-False-Claims-Act-Liability}{Wiley} \href{https://pmc.ncbi.nlm.nih.gov/articles/PMC8247552/}{PubMed Central} but prosecution remains rare and targeted at the most egregious financial fraud, not systemic methodological failure. \href{https://www.hugheshubbard.com/news/second-circuit-ratchets-up-false-claims-act-exposure-for-grant-recipients}{Hugheshubbard}

The Trump administration's 2025--2026 NIH actions --- proposed 40\% budget cuts, grant terminations, indirect cost caps --- were driven by ideology (targeting DEI and vaccine research) and cost-cutting, not by research integrity concerns. RFK Jr.~cited the amyloid fraud during his confirmation hearing but proposed no constructive reforms. \href{https://alzimpact.org/Contrary-to-Clear-Verifiable-Facts-Secretary-Kennedy-Continues-to-Repeat-Harmful-Myths}{Alzimpact} \href{https://www.theburningplatform.com/2026/03/26/the-great-alzheimers-scam-and-the-proven-cures-theyve-buried-for-billions-2/}{The Burning Platform} The scientific community's response --- ``Stand Up for Science'' rallies, court challenges --- focused on defending funding levels rather than addressing the structural dysfunction those funds perpetuate. A federal court ruled the administration \textbf{illegally froze \$8 billion} in grants, \href{https://www.govexec.com/oversight/2025/08/trump-illegally-froze-1800-nih-medical-research-grants-congress-watchdog-says/407296/}{GovExec} and Congress ultimately set the FY2026 NIH budget at \$48.7 billion.

\subsection{Why this never becomes a movement}\label{19_200_billion-why-this-never-becomes-a-movement}

The deepest puzzle --- why documented fraud affecting millions of patients and consuming billions in public funds generates no sustained political response --- has a multi-layered answer that goes beyond simple institutional capture.

First, the issue has \textbf{no natural political home}. Conservatives tend to want to cut science funding entirely; liberals tend to defend the scientific establishment uncritically. Research accountability doesn't map onto existing coalitions. Post-COVID, ``defending science'' became a politically charged identity marker, making constructive criticism risky. \href{https://issues.org/new-politics-science-mills-st-clair/}{Issues} RFK Jr.'s weaponization of legitimate integrity concerns alongside vaccine misinformation has, as multiple scientists noted, ``poisoned the well'' for reform --- anyone questioning research integrity risks being lumped with anti-science movements.

Second, the scientific establishment successfully frames each fraud case as \textbf{``the system working.''} The fraud was eventually detected; papers were eventually retracted; the researcher eventually departed. A \emph{Brain Communications} editorial called the Lesné case ``an illustration of the scientific process working towards an accurate picture of the world'' \href{https://pmc.ncbi.nlm.nih.gov/articles/PMC9445174/}{PubMed Central} --- ignoring that detection took 16 years, came from outside the system, and produced no structural change.

Third, the victims are diffuse and the causal chain is long. There is no identifiable patient who can point to the Lesné paper and say ``this is why I wasn't cured.'' Alzheimer's patients and their families direct their energy toward hope --- more funding, more research --- not toward questioning whether existing research was valid. \href{https://www.barnesandnoble.com/w/doctored-charles-piller/1145681914}{Barnes \& Noble} The Alzheimer's Association reinforces this by attacking critics and maintaining that ``less than 14\%'' of NIH Alzheimer's projects focused on amyloid as a therapeutic target, \href{https://alzimpact.org/Contrary-to-Clear-Verifiable-Facts-Secretary-Kennedy-Continues-to-Repeat-Harmful-Myths}{Alzimpact} a framing disputed by independent analysis.

Fourth, the \textbf{concentrated beneficiaries are extraordinarily well-organized}. The biomedical research enterprise generates approximately \$200+ billion annually across NIH grants (\$48B), university overhead (\$9B+), scientific publishing (\$19--28B globally), laboratory equipment (Thermo Fisher alone: \$43B revenue), \href{https://en.wikipedia.org/wiki/Thermo_Fisher_Scientific}{Wikipedia} animal breeding (Charles River: \$4B revenue), and downstream pharmaceutical R\&D. Each component has professional associations, lobbying operations, and direct connections to every Congressional district. Against this stands no organized constituency for ``research that actually produces cures.''

\subsection{Conclusion}\label{19_200_billion-conclusion}

The biomedical research system represents perhaps the most successful instance of institutional capture in modern governance --- a \$200+ billion enterprise that has converted the public's desperate desire for medical cures into a self-sustaining funding ecosystem where \textbf{the primary output is career advancement for researchers, revenue for institutions, and profits for publishers and suppliers}, with clinical translation as an occasional byproduct. The evidence is not ambiguous: 95\% translation failure rates, \href{https://journals.plos.org/plosbiology/article?id=10.1371/journal.pbio.3002667}{PLOS} 11\% replication rates for landmark studies, \href{https://issues.org/ending-reproducibility-crisis-medical-research-brownlee-bielekova/}{Issues} 72\% of researchers acknowledging a crisis, \href{https://www.nature.com/articles/d41586-024-04253-w}{Nature} \href{https://journals.plos.org/plosbiology/article?id=10.1371/journal.pbio.3002870}{PLOS} zero approved treatments for methamphetamine addiction, \href{https://pubmed.ncbi.nlm.nih.gov/34344291/}{PubMed} a 99.6\% failure rate in Alzheimer's drug development, \href{https://www.scientificamerican.com/article/why-alzheimer-s-drugs-keep-failing/}{Scientific American} and a documented fraud that consumed billions with no consequences for anyone.

The system persists because it has no natural antagonist. Every entity that should provide accountability --- universities, ORI, journals, Congress, patient groups, media --- is either financially dependent on the system or structurally incapable of challenging it. The reform ecosystem (HHMI, Fast Grants, ARPA-H, FROs, registered reports) demonstrates that alternative models produce better science, but these remain peripheral experiments operating at a fraction of NIH's scale. The most promising structural development --- ARPA-H's DARPA-inspired model --- has a \$1.5 billion budget \href{https://www.science.org/content/article/final-nih-budget-2024-essentially-flat}{Science} against NIH's \$48 billion.

What distinguishes this from other regulatory capture scenarios is the moral dimension: the captured institution's stated purpose is curing disease in human beings. Every year the system persists in its current form represents not just wasted money but unmade discoveries --- treatments for Alzheimer's, addiction, depression, and cancer that a differently structured system might have produced. The question is not whether anyone has identified the problem. They have \href{https://retractionwatch.com/2017/04/04/failure-essential-part-science-qa-author-new-book-reproducibility/}{Retraction Watch} --- Ioannidis, Piller, Harris, Buck, Collison, and others have documented it exhaustively. The question is whether any political force exists with both the motivation and the power to restructure \$200 billion in entrenched interests. As of 2026, the honest answer is no.
